\documentclass{article}
\usepackage[utf8]{vietnam}
\usepackage{times}
\usepackage{amsmath,amssymb}
\usepackage[version=4]{mhchem}
\usepackage{tabularx}
\usepackage{array}
\usepackage{booktabs}
\usepackage{xcolor}
\usepackage[
  paperwidth=29.7cm,
  paperheight=16.706cm,
  top=1.0cm,
  bottom=1.0cm,
  left=1.2cm,
  right=1.2cm,
  footskip=0pt,
  headheight=0pt,
  headsep=0pt
]{geometry}

\pagestyle{empty}
\setlength{\parindent}{0pt}
\setlength{\parskip}{0pt}

\definecolor{ansRed}{RGB}{198, 40, 40}

\begin{document}

\fontsize{26}{34}\selectfont

\textbf{Câu 1.} Cho dãy các kí hiệu các nguyên tố sau: He, Ca, K, Fe, Cu, S. Theo thứ tự tên của các nguyên tố lần lượt là:
\vspace{6pt}

\begin{tabularx}{\linewidth}{@{}X X@{}}
\textbf{A.} Helium, Caesium, Potassium, Copper, Iron, Silicon. & \textbf{B.} Helium, Calcium, Potassium, Iron, Copper, Sulfur.\\[8pt]
\textbf{C.} Helium, Calcium, Krypton, Iron, Copper, Silicon. & \textbf{D.} Hydogen, Caesium, Krypton, Iron, Cobalt, Sulfur.
\end{tabularx}

\vspace*{\fill}

\textbf{Câu 2.} Các nguyên tố trong cùng 1 chu kì thì có cùng:
\vspace{6pt}

\begin{tabularx}{\linewidth}{@{}X X@{}}
\textbf{A.} Số electron lớp ngoài cùng. & \textbf{B.} Khối lượng nguyên tử\\[8pt]
\textbf{C.} Số điện tích hạt nhân. & \textbf{D.} Số lớp electron.
\end{tabularx}

\newpage

\fontsize{26}{34}\selectfont

\textbf{Câu 3.} Đơn chất là những chất được tạo nên từ bao nhiêu nguyên tố hóa học?
\vspace{6pt}

\begin{tabularx}{\linewidth}{@{}X X@{}}
\textbf{A.} Từ 2 nguyên tố. & \textbf{B.} Từ 3 nguyên tố.\\[8pt]
\textbf{C.} Từ 4 nguyên tố trở lên. & \textbf{D.} Từ 1 nguyên tố.
\end{tabularx}

\vspace*{\fill}

\textbf{Câu 4.} Hạt nhân một nguyên tử fluorine có 9 proton và 10 neutron. Khối lượng của nguyên tử fluorine xấp xỉ bằng:
\vspace{6pt}

\begin{tabularx}{\linewidth}{@{}X X@{}}
\textbf{A.} 9 amu & \textbf{B.} 10 amu.\\[8pt]
\textbf{C.} 19 amu & \textbf{D.} 28 amu.
\end{tabularx}

\newpage

\fontsize{26}{34}\selectfont

\textbf{Câu 5.} Một bình khí oxygen chứa:
\vspace{6pt}

\begin{tabularx}{\linewidth}{@{}X X@{}}
\textbf{A.} Các phân tử O$_{2}$. & \textbf{B.} Các nguyên tử oxygen riêng rẽ không liên kết với nhau.\\[8pt]
\textbf{C.} Một phân tử O$_{2}$ & \textbf{D.} Một đại phân tử khổng lồ chứa rất nhiều nguyên tử oxygen.
\end{tabularx}

\vspace*{\fill}

\textbf{Câu 6.} Loại hạt nào sau đây đặc trưng cho 1 nguyên tố hóa học?
\vspace{6pt}

\begin{tabularx}{\linewidth}{@{}X X@{}}
\textbf{A.} hạt electron & \textbf{B.} hạt proton\\[8pt]
\textbf{C.} Hạt electron và proton & \textbf{D.} hạt electron và neutron
\end{tabularx}

\newpage

\fontsize{26}{34}\selectfont

\textbf{Câu 7.} Một nguyên tử carbon có khối lượng là 1,9926 $\times$  10$^{-23}$ gam. Khối lượng nguyên tử carbon tính theo đơn vị amu là: (biết 1 amu = 1,6605 $\times$  10$^{-24}$ gam).
\vspace{6pt}

\begin{tabularx}{\linewidth}{@{}X X@{}}
\textbf{A.} 12 amu. & \textbf{B.} 24 amu.\\[8pt]
\textbf{C.} 6 amu. & \textbf{D.} 48 amu.
\end{tabularx}

\vspace*{\fill}

\textbf{Câu 8.} Thời tiết nồm ẩm vào cuối mùa xuân có bản chất là hiện tượng:
\vspace{6pt}

\begin{tabularx}{\linewidth}{@{}X X@{}}
\textbf{A.} ngưng tụ hơi nước. & \textbf{B.} bay hơi nước ở thể lỏng.\\[8pt]
\textbf{C.} đông đặc nước ở thể lỏng. & \textbf{D.} nóng chảy nước ở thể rắn.
\end{tabularx}

\newpage

\fontsize{26}{34}\selectfont

\textbf{Câu 9.} Chất nào sau đây không phải là hợp chất cộng hóa trị?
\vspace{6pt}

\begin{tabularx}{\linewidth}{@{}X X@{}}
\textbf{A.} SO$_{2}$. & \textbf{B.} CO$_{2}$.\\[8pt]
\textbf{C.} NaCl. & \textbf{D.} H$_{2}$O.
\end{tabularx}

\vspace*{\fill}

\textbf{Câu 10.} Cho biết tổng số hạt trong hạt nhân nguyên tử Nitrogen là 14, số đơn vị điện tích hạt nhân là 7. Số hạt neutron và electron của nguyên tử Nitrogen:
\vspace{6pt}

\begin{tabularx}{\linewidth}{@{}X X@{}}
\textbf{A.} 7 và 14 & \textbf{B.} 14 và 7\\[8pt]
\textbf{C.} 7 và 7. & \textbf{D.} 14 và 14.
\end{tabularx}

\newpage

\fontsize{26}{34}\selectfont

\textbf{Câu 11.} Cho các phát biểu:\\ (I) Nguyên tử trung hoà về điện.\\ (II) Khối lượng của nguyên tử tập trung chủ yếu ở hạt nhân.\\ (III) Trong nguyên tử, số hạt mang điện tích dương bằng số hạt mang điện tích âm nên số hạt electron bằng số hạt neutron.\\ (IV) Vỏ nguyên tử, gồm các lớp electron có khoảng cách khác nhau đối với hạt nhân.\\ Trong các phát biểu trên, số phát biểu đúng là:
\vspace{6pt}

\begin{tabularx}{\linewidth}{@{}X X@{}}
\textbf{A.} 1. & \textbf{B.} 2.\\[8pt]
\textbf{C.} 3. & \textbf{D.} 4.
\end{tabularx}

\vspace*{\fill}

\textbf{Câu 12.} Cho các hợp chất sau: MgCl$_{2}$, Na$_{2}$O, NH$_{3}$, HCl, KCl. Hợp chất nào sau đây có liên kết cộng hóa trị?
\vspace{6pt}

\begin{tabularx}{\linewidth}{@{}X X@{}}
\textbf{A.} NH$_{3}$ và HCl & \textbf{B.} Na$_{2}$O và NH$_{3}$.\\[8pt]
\textbf{C.} MgCl$_{2}$ và Na$_{2}$O. & \textbf{D.} HCl và KCl
\end{tabularx}

\newpage

\fontsize{26}{34}\selectfont

\textbf{Câu 13.} Hợp chất Na$_{x}$SO$_{4}$ có khối lượng phân tử là 142 amu. Giá trị của x là:
\vspace{6pt}

\begin{tabularx}{\linewidth}{@{}X X@{}}
\textbf{A.} 1. & \textbf{B.} 2.\\[8pt]
\textbf{C.} 3. & \textbf{D.} 4.
\end{tabularx}

\vspace*{\fill}

\textbf{Câu 14.} Nguyên tố X có 12 proton, 3 lớp electron, có 2 e lớp ngoài cùng. Vị trí của nguyên tố X trong bảng tuần hoàn là:
\vspace{6pt}

\begin{tabularx}{\linewidth}{@{}X X@{}}
\textbf{A.} X ở ô số 12, chu kì 3, nhóm III & \textbf{B.} X ở ô số 12, chu kì 2, nhóm I.\\[8pt]
\textbf{C.} X ở ô số 12, chu kì 3, nhóm II & \textbf{D.} X ở ô số 12, chu kì 2, nhóm III
\end{tabularx}

\newpage

\fontsize{26}{34}\selectfont

\textbf{Câu 15.} Trong phân tử nước, cứ 16,0 g oxygen có tương ứng 2,0 g hydrogen. Một giọt nước chứa 0,1 g hydrogen thì khối lượng của oxygen có trong giọt nước đó là:
\vspace{6pt}

\begin{tabularx}{\linewidth}{@{}X X@{}}
\textbf{A.} 1,6 g. & \textbf{B.} 1,2 g.\\[8pt]
\textbf{C.} 0,9g. & \textbf{D.} 0,8 g.
\end{tabularx}

\vspace*{\fill}

\textbf{Câu 16.} Khối lượng phân tử sulfur dioxide (SO$_{2}$) là:
\vspace{6pt}

\begin{tabularx}{\linewidth}{@{}X X@{}}
\textbf{A.} 64 amu & \textbf{B.} 48 amu\\[8pt]
\textbf{C.} 16 amu & \textbf{D.} 80 amu
\end{tabularx}

\newpage

\fontsize{26}{34}\selectfont
\textbf{Câu 17 (Đúng/Sai).} Một số đặc điểm của đơn chất:
\vspace{14pt}

\begin{tabularx}{\linewidth}{@{}p{0.88\linewidth} c@{}}
\textbf{a)} Các nguyên tử trong đơn chất kim loại sắp xếp khít nhau theo một trật tự xác định và có vai trò như phân tử. & $[\quad]$\\[12pt]\textbf{b)} Các đơn chất kim loại có tính chất vật lí chung như: Dẫn điện, dẫn nhiệt, có ánh kim,\dots  & $[\quad]$\\[12pt]\textbf{c)} Đơn chất khí hiếm tồn tại trong tự nhiên ở dạng phân tử. & $[\quad]$\\[12pt]\textbf{d)} Các đơn chất phi kim có tính chất khác với các đơn chất kim loại. & $[\quad]$
\end{tabularx}

\newpage

\fontsize{26}{34}\selectfont
\textbf{Câu 18 (Đúng/Sai).} Nếu một nguyên tử có 14 electron thì:
\vspace{14pt}

\begin{tabularx}{\linewidth}{@{}p{0.88\linewidth} c@{}}
\textbf{a)} Nguyên tử đó có 14 proton. & $[\quad]$\\[12pt]\textbf{b)} Nguyên tử đó có 2 lớp electron. & $[\quad]$\\[12pt]\textbf{c)} Nguyên tử có số đơn vị điện tích hạt nhân là 14. & $[\quad]$\\[12pt]\textbf{d)} Nguyên tử có số electron lớp ngoài cùng là 3. & $[\quad]$
\end{tabularx}

\newpage

\fontsize{26}{34}\selectfont
\textbf{Câu 19 (Đúng/Sai).} Cấu tạo nguyên tử
\vspace{14pt}

\begin{tabularx}{\linewidth}{@{}p{0.88\linewidth} c@{}}
\textbf{a)} Nguyên tử gồm 3 loại hạt cơ bản là proton, neutron và electron. & $[\quad]$\\[12pt]\textbf{b)} Lớp vỏ nguyên tử gồm 2 loại hạt là electron và neutron. & $[\quad]$\\[12pt]\textbf{c)} Hạt nhân nguyên tử gồm 2 loại hạt là proton và neutron. & $[\quad]$\\[12pt]\textbf{d)} Nguyên tử luôn trung hòa về điện nên số hạt proton luôn bằng số hạt neutron. & $[\quad]$
\end{tabularx}

\newpage

\fontsize{30}{40}\selectfont
\textbf{Câu 20 (Trả lời ngắn).} Cho hợp chất P$_{2}$O$_{5}$. Khối lượng phân tử của hợp chất trên là bao nhiêu amu?
\vspace{22pt}

\textbf{Trả lời:} $\dots\dots\dots\dots\dots\dots\dots\dots\dots\dots$

\newpage

\fontsize{30}{40}\selectfont
\textbf{Câu 21 (Trả lời ngắn).} Thành phần phần trăm theo khối lượng của nguyên tố P trong hợp chất P$_{2}$O$_{5}$ là bao nhiêu phần trăm ? (Tính ra kết quả số thập phân và làm tròn đến 2 chữ số sau dấu phẩy ).
\vspace{22pt}

\textbf{Trả lời:} $\dots\dots\dots\dots\dots\dots\dots\dots\dots\dots$

\newpage

\fontsize{30}{40}\selectfont
\textbf{Câu 22 (Trả lời ngắn).} Nguyên tử Aluminium có điện tích hạt nhân là 13+. Trong nguyên tử, số hạt mang điện nhiều hơn số hạt không mang điện là 12 hạt. Số hạt neutron của nguyên tử Aluminium là bao nhiêu ?
\vspace{22pt}

\textbf{Trả lời:} $\dots\dots\dots\dots\dots\dots\dots\dots\dots\dots$

\newpage

\fontsize{30}{40}\selectfont
\textbf{Câu 23 (Trả lời ngắn).} Trong hạt nhân nguyên tử fluorine có 9 proton. Số electron ở lớp ngoài cùng của vỏ nguyên tử fluorine là bao nhiêu?
\vspace{22pt}

\textbf{Trả lời:} $\dots\dots\dots\dots\dots\dots\dots\dots\dots\dots$

\newpage

\fontsize{30}{40}\selectfont
\textbf{Câu 24 (Trả lời ngắn).} Cho hợp chất K$_{2}$S. Khối lượng phân tử của hợp chất trên là bao nhiêu amu?
\vspace{22pt}

\textbf{Trả lời:} $\dots\dots\dots\dots\dots\dots\dots\dots\dots\dots$

\newpage

\fontsize{30}{40}\selectfont
\textbf{Câu 25 (Trả lời ngắn).} Một nguyên tử có 10 proton trong hạt nhân. Theo mô hình nguyên tử của Rơ- dơ- pho- Bo, nguyên tử trên có bao nhiêu lớp electron ?
\vspace{22pt}

\textbf{Trả lời:} $\dots\dots\dots\dots\dots\dots\dots\dots\dots\dots$

\newpage

\fontsize{18}{24}\selectfont
\begin{center}
{\fontsize{24}{30}\selectfont\textbf{BẢNG ĐÁP ÁN CHÍNH THỨC -- ĐỀ KIỂM TRA GIỮA KÌ I (MÃ ĐỀ 01)}}\\[10pt]
\end{center}

\textbf{1. PHẦN I: CÂU TRẮC NGHIỆM NHIỀU PHƯƠNG ÁN LỰA CHỌN}\\[4pt]
\renewcommand{\arraystretch}{1.2}
\begin{tabularx}{\linewidth}{|*{16}{>{\centering\arraybackslash}X|}}
\hline
\textbf{1} & \textbf{2} & \textbf{3} & \textbf{4} & \textbf{5} & \textbf{6} & \textbf{7} & \textbf{8} & \textbf{9} & \textbf{10} & \textbf{11} & \textbf{12} & \textbf{13} & \textbf{14} & \textbf{15} & \textbf{16} \\
\hline
\color{ansRed}\textbf{B} & \color{ansRed}\textbf{D} & \color{ansRed}\textbf{D} & \color{ansRed}\textbf{C} & \color{ansRed}\textbf{A} & \color{ansRed}\textbf{B} & \color{ansRed}\textbf{A} & \color{ansRed}\textbf{A} & \color{ansRed}\textbf{C} & \color{ansRed}\textbf{C} & \color{ansRed}\textbf{C} & \color{ansRed}\textbf{A} & \color{ansRed}\textbf{B} & \color{ansRed}\textbf{C} & \color{ansRed}\textbf{D} & \color{ansRed}\textbf{A} \\
\hline
\end{tabularx}

\vspace{14pt}

\textbf{2. PHẦN II: CÂU TRẮC NGHIỆM ĐÚNG SAI}\\[6pt]
\begin{tabularx}{\linewidth}{@{}*{3}{>{\raggedright\arraybackslash}X}@{}}
\textbf{Câu 17:} [\color{ansRed}\textbf{Đ -- Đ -- S -- Đ}\color{black}] &
\textbf{Câu 18:} [\color{ansRed}\textbf{Đ -- S -- Đ -- S}\color{black}] &
\textbf{Câu 19:} [\color{ansRed}\textbf{Đ -- S -- Đ -- S}\color{black}]
\end{tabularx}

\vspace{14pt}

\textbf{3. PHẦN III: CÂU TRẮC NGHIỆM TRẢ LỜI NGẮN}\\[6pt]
\begin{tabularx}{\linewidth}{@{}*{6}{>{\centering\arraybackslash}X}@{}}
\hline
\textbf{Câu 20} & \textbf{Câu 21} & \textbf{Câu 22} & \textbf{Câu 23} & \textbf{Câu 24} & \textbf{Câu 25} \\
\hline
\color{ansRed}\textbf{142} & \color{ansRed}\textbf{43,66\%} & \color{ansRed}\textbf{14} & \color{ansRed}\textbf{7} & \color{ansRed}\textbf{110} & \color{ansRed}\textbf{2} \\
\hline
\end{tabularx}

\end{document}
